\documentclass[12pt]{article}
\usepackage{fancyhdr}
\pagestyle{fancy}
\fancyhf{}
\fancyfoot[C]{\thepage}
\fancyfoot[L]{\scriptsize CC BY-NC-SA 4.0 \textbullet\ Leo Liang}
\renewcommand{\headrulewidth}{0pt}
\usepackage[margin=1in]{geometry}
\usepackage{amsmath, amssymb}
\usepackage{array, booktabs}
\usepackage{pgffor}
\usepackage{graphicx}
\parindent=0pt
\parskip=1em
\newcolumntype{C}[1]{>{\centering\arraybackslash}p{#1}}
% writing lines
\newcommand{\wl}{\par\vspace{5mm}\noindent\rule{\linewidth}{0.4pt}\par\vspace{1mm}}
\newcommand{\wls}[1]{\foreach \i in {1,...,#1}{\wl}}
% ─────────────────────────────────────────────────────────────────────────────
\begin{document}
\pagestyle{fancy}
\begin{center}
  {\Large\bfseries Homework 2 --- The Language of Sets}\\[4pt]
  Week 2\\[6pt]
  Name:\;\underline{\hspace{6.5cm}}\quad
\end{center}
\hrule\bigskip
% =============================================================================
\section*{Instructions}
Use \textbf{pen}. Write questions and surprises in the margin in a different color.
Some things in the reading are intentionally imprecise --- find them. This reading is difficult,
if there are topics you don't understand, you should use the internet or AI to help. AI can be a good teacher
if you ask it specific questions. \textbf{DO NOT} use AI
to solve the problems without trying yourself first.
% =============================================================================
\section{Readings}
Take notes in your notebook on important definitions and examples that you find helpful. Taking notes
is a skill in itself. As you read the following content, ask questions in your head (using what you learned from the \textbf{question
compass}) and try answering it. If you can't answer your own question, re-read the text. If you still can't
answer it after re-reading, write your question down in a \textbf{different color pen} in the margins and we will go over in it class.
% =============================================================================
% DROP-IN SECTION: Maximal Interpretation
% =============================================================================
\subsection*{Maximal Interpretation}
\noindent\textbf{Every Clue Earns Its Place}

In class we looked at a problem about two Minecraft bedrooms.
One student understood the problem so that Bedroom 2's width (8 blocks) was also the
length of the shared wall --- a perfectly valid logical interpretation.
But in that reading, the last sentence of the problem
(``you haven't decided how wide Bedroom 1 will be yet'')
was doing nothing. No variable was needed. No unknown existed.

That interpretation wasn't \textit{wrong} --- it just wasn't \textbf{maximal}.

\medskip
\noindent\textbf{\textit{Definition.}}\quad An interpretation is \textbf{maximal}
if every piece of given information is used.
An interpretation that leaves a clue unused is called \textbf{wasteful}.

\medskip
In a well-written problem, every sentence is there for a reason.
Before you commit to a reading, ask yourself:
\textit{Is every piece of given information doing something in my interpretation?}
If a sentence is just sitting there, look for a better reading.

\medskip
There is a hierarchy of how badly an interpretation can fail:

\medskip
\renewcommand{\arraystretch}{1.6}
\begin{tabular}{@{} p{3.5cm} p{9.5cm} @{}}
\toprule
\textbf{Type} & \textbf{What it means} \\
\midrule
Contradiction   & Your interpretation makes two clues logically impossible
                  at the same time. Ruled out entirely. \\
Wasteful        & Your interpretation is consistent, but leaves at least one
                  clue unused. Probably wrong --- keep looking. \\
Maximal         & Every clue is active. This is the interpretation to work with. \\
\bottomrule
\end{tabular}
\renewcommand{\arraystretch}{1}

\medskip
\textit{Note:} A maximal interpretation is not always unique.
If two different readings each use all the given information,
both deserve consideration --- and the problem itself may be ambiguous.
That is worth flagging too. Often times it is just a poorly written question.

\bigskip
\hrule
\subsection*{Interpretation Practice}

\textbf{I1.}\quad
Write a problem or statement such that there exists a maximal interpretation and wasteful interpretation (can be more than two). Explain the possible interpretations
and what makes them maximal versus wasteful.
\wls{3}

\subsection*{What Is a Set?}
\textit{\textbf{Definition.}}\quad A \textbf{set} is an unordered collection of \textbf{distinct} objects.
Each object is called an \textbf{element} (or \textbf{member}) of the set.
\[
  a \in A \quad \text{means} \quad \text{``}a\text{ is an element of }A\text{''}
\]
\[
  a \notin A \quad \text{means} \quad \text{``}a\text{ is NOT an element of }A\text{''}
\]
\textbf{Example.} Let $S = \{1, 3, 5\}$. Then $3 \in S$, but $4 \notin S$.

One important fact: a set has \textbf{no repeated elements} and \textbf{no order}.
The sets $\{1,2,3\}$, $\{3,1,2\}$, and $\{1,1,2,3\}$ are all the same set.

\bigskip
\subsection*{Two Ways to Write a Set}
\textbf{Roster notation} --- list the elements explicitly:
\[
  A = \{2, 4, 6, 8\}
\]
\textbf{Set-builder notation} --- describe the property elements must satisfy:
\[
  A = \{x \in \mathbb{Z} \mid x \text{ is even and } 0 < x \leq 8\}
\]
Read the vertical bar $\mid$ as \textbf{``such that.''} The full phrase:
\textit{``the set of all $x$ in $\mathbb{Z}$ such that $x$ is even and between 0 and 8.''}

\textbf{Example.}\quad $\{x \in \mathbb{R} \mid x^2 = 9\} = \{-3,\; 3\}$.
Why are there two elements? Because both $(-3)^2 = 9$ and $3^2 = 9$.

\newpage
\subsection*{The Number Sets You Will See Constantly}
\medskip
\renewcommand{\arraystretch}{1.6}
\begin{tabular}{@{} p{1.4cm} p{4.2cm} p{8cm} @{}}
\toprule
Symbol & Name & Description \\
\midrule
$\mathbb{N}$ & Natural numbers    & $\{1,\, 2,\, 3,\, 4,\, \ldots\}$ \\
$\mathbb{Z}$ & Integers           & $\{\ldots,\, {-2},\, {-1},\, 0,\, 1,\, 2,\, \ldots\}$ \\
$\mathbb{Q}$ & Rational numbers   & All fractions $\tfrac{p}{q}$ with $p,q \in \mathbb{Z}$, $q \neq 0$ \\
$\mathbb{R}$ & Real numbers       & Every point on the number line \\
$\emptyset$  & The empty set      & $\{\;\}$ --- a perfectly valid set containing nothing \\
\bottomrule
\end{tabular}
\renewcommand{\arraystretch}{1}

\medskip
These sets are \textbf{nested}:
$\mathbb{N} \subseteq \mathbb{Z} \subseteq \mathbb{Q} \subseteq \mathbb{R}$.
Every natural number is an integer, every integer is rational, and so on.

\medskip
\begin{center}
    \textit{[Number-classification diagram omitted from the public release.]}
\end{center}

\bigskip
\newpage
\subsection*{Subsets}
Set $A$ is a \textbf{subset} of set $B$, written $A \subseteq B$,
if every element of $A$ is also in $B$.

Formally: $A \subseteq B$ means: for all $x$, if $x \in A$ then $x \in B$.

\textbf{Examples:}
\begin{itemize}
  \item $\{2,4\} \subseteq \{1,2,3,4,5\}$? \textbf{Yes.} Both 2 and 4 appear.
  \item $\{2,7\} \subseteq \{1,2,3,4,5\}$? \textbf{No.} $7 \notin \{1,2,3,4,5\}$.
  \item $\{1,2,3\} \subseteq \{1,2,3\}$? \textbf{Yes}. they are equal in this case.
  \item $\emptyset \subseteq \{1,2,3\}$? \textbf{Yes}. All sets have nothing in it.
\end{itemize}

\subsection*{Set Operations}
Given sets $A$ and $B$:

\medskip
\renewcommand{\arraystretch}{1.8}
\begin{tabular}{@{} p{3.5cm} p{3cm} p{7cm} @{}}
\toprule
Operation & Notation & Meaning \\
\midrule
Union        & $A \cup B$ & Everything in $A$ OR $B$ (or both) \\
Intersection & $A \cap B$ & Everything in BOTH $A$ AND $B$ \\
Difference   & $A \setminus B$ & In $A$ but NOT in $B$ \\
Complement   & $A^c$      & Everything in $U$ that is NOT in $A$ \\
\bottomrule
\end{tabular}
\renewcommand{\arraystretch}{1}

\textit{\textbf{Definition.}}\quad A \textbf{universal set} $U$ is a declared set that contains every object
under consideration. All other sets in the problem satisfy $A \subseteq U$.

\medskip
\textbf{Example.}\quad If we are working with integers from 1 to 10,
we write $U = \{1, 2, 3, \ldots, 10\}$.
Any set we discuss --- say $A = \{2, 4, 6\}$ --- must be a subset of $U$.

\medskip
\textbf{Why it matters.}\quad Union, intersection, and set difference
depend only on $A$ and $B$. But complement depends on $U$:
\[
  A^c = \{x \in U \mid x \notin A\}
\]
Change $U$, and $A^c$ changes. Without declaring $U$, the complement is undefined.

\medskip
\textbf{Example.}\quad Let $U = \{1,2,3,4,5\}$, $A = \{1,2,3\}$, $B = \{2,3,4\}$.
\[
  A \cup B = \{1,2,3,4\},\qquad
  A \cap B = \{2,3\},\qquad
  A \setminus B = \{1\},\qquad
  A^c = \{4,5\}
\]

\medskip
\textbf{\textit{Definition}.}\quad The \textbf{cardinality} of a set $A$,
written $|A|$, is the number of elements in $A$.

Using the example above, fill in each cardinality:

\medskip
\renewcommand{\arraystretch}{2.2}
\begin{tabular}{|C{3.5cm}|C{3.5cm}|C{3.5cm}|C{3.5cm}|}
\hline
$|A|$ & $|B|$ & $|A \cup B|$ & $|A \cap B|$ \\
\hline
& & & \\
\hline
\end{tabular}
\renewcommand{\arraystretch}{1}

\bigskip
\subsection*{Cartesian Products}
\textbf{\textit{Definition.}}\quad The \textbf{Cartesian product} $A \times B$ (read ``$A$ cross $B$'') is
the set of ALL ordered pairs $(a,b)$ where $a \in A$ and $b \in B$:
\[
  A \times B = \{(a,b) \mid a \in A \text{ and } b \in B\}
\]
\textbf{Example.}\quad $A = \{1,2\}$, $B = \{\text{red},\, \text{blue},\, \text{green}\}$:
\[
  A \times B = \{(1,\text{red}),\,(1,\text{blue}),\,(1,\text{green}),\,
                 (2,\text{red}),\,(2,\text{blue}),\,(2,\text{green})\}
\]
So $|A \times B| = 2 \times 3 = 6$.

\medskip
\textbf{You already know one Cartesian product:} the coordinate plane is
$\mathbb{R} \times \mathbb{R} = \mathbb{R}^2$.
Every point $(x, y)$ is an ordered pair --- one from each copy of $\mathbb{R}$.

\medskip
\textbf{Big preview.}\quad In the next worksheet, we define a
\textit{function} $f\colon A \to B$ as a special subset of $A \times B$ ---
a set of ordered pairs where each $a \in A$ appears exactly once on the left.
Keep that in mind. By the end of next class, this sentence will make complete sense.

\subsection*{Why You Cannot Just Add}
Suppose 13 students like Math and 9 students like Science.
Can you conclude that 22 students like at least one subject?

Not necessarily --- some students might like \textit{both}.
If 5 students like both, then adding $13 + 9 = 22$ counts those
5 students twice: once in the Math group and once in the Science group. To fix the overcount, subtract the overlap once. In this situation, there are $13 + 9 - 5 = 17$ students that like at least one subject.
Notice mathematically what we did was:
\[
  |\text{Math} \cup \text{Science}|
  \;=\;
  |\text{Math}| + |\text{Science}| - |\text{Math} \cap \text{Science}|
  \;=\;
  13 + 9 - 5
  \;=\;
  17
\]
Also notice that the phrase ``students that like \textbf{at least} one subject'' translates to \[|\text{Math} \cup \text{Science}|.\]

\subsubsection*{The Inclusion-Exclusion Formula}
The previous problem demonstrates an important formula in set theory, the \textbf{inclusion-exclusion formula}.
This works for any two sets:
\[
  |A \cup B| = |A| + |B| - |A \cap B|
\]
\textbf{Why it works.}\quad Every element of $A \cap B$
gets added in $|A|$, added again in $|B|$, then subtracted once.
Net result: counted exactly once.
Every element in only $A$ or only $B$ is never subtracted.
Net result: also counted exactly once.

\medskip
\textbf{Example.}\quad
On a Minecraft server, 13 players have built a house and 9 have
defeated the Ender Dragon. 5 players have done both.
How many players have done at least one?
\[
  |H \cup D| = 13 + 9 - 5 = 17
\]
If there are 25 players total, then $25 - 17 = 8$ have done neither.

\medskip
\textbf{Practice.}\quad
At a summer camp, 24 campers went hiking and 19 went swimming.
8 campers did both. There are 40 campers total.

(a)\quad How many campers did at least one activity?
\wls{2}

(b)\quad How many campers did neither?
\wl

\vspace{4cm}
\newpage

% =============================================================================
\subsection*{Foundation Practice}

\textbf{1.}\quad
Membership, Notation, and Set-Builder Form.

\textbf{(a)}\quad Let $S = \{3,\, 6,\, 9,\, 12,\, 15\}$.
Fill in $\in$ or $\notin$:

\medskip
\renewcommand{\arraystretch}{2.2}
\begin{tabular}{|C{3.5cm}|C{3.5cm}|C{3.5cm}|C{3.5cm}|}
\hline
$6 \mathrel{\underline{\hspace{0.8cm}}} S$ &
$0 \mathrel{\underline{\hspace{0.8cm}}} S$ &
$9 \mathrel{\underline{\hspace{0.8cm}}} S$ &
$18 \mathrel{\underline{\hspace{0.8cm}}} S$ \\
\hline
\end{tabular}
\renewcommand{\arraystretch}{1}

\medskip
\textbf{(b)}\quad Write $S$ from (a) in set-builder notation using $\mathbb{N}$.
(Hint: what property do ALL elements of $S$ share?)
\wls{2}

\textbf{(c)}\quad Write each in roster notation.

\medskip
\hspace{1.5em}(i)\quad $\{x \in \mathbb{Z} \mid -2 \leq x \leq 2\}$
\hfill $= \;\underline{\hspace{6cm}}$

\medskip
\hspace{1.5em}(ii)\quad $\{x \in \mathbb{Z} \mid x^2 < 10\}$
\hfill $= \;\underline{\hspace{6cm}}$

\medskip
\hspace{1.5em}(iii)\quad $\{x \in \mathbb{N} \mid x \text{ divides } 12\}$
\hfill $= \;\underline{\hspace{6cm}}$

\textbf{(d)}\quad Write $\{x \in \mathbb{R} \mid x^2 = 4\}$ in roster notation.
\[
  \{x \in \mathbb{R} \mid x^2 = 4\} \;=\; \underline{\hspace{5cm}}
\]

\bigskip
\textbf{2.}\quad Subsets.

\textbf{(a)}\quad List ALL subsets of $\{p,\, q,\, r\}$. Include the empty set.
\wls{3}

How many subsets are there total? \;\hrulefill

If you see a pattern, try to state it in terms of the number of elements in the original set.
\wl

\textbf{(b)}\quad Answer T\,/\,F and justify in one sentence.

\medskip
T\,/\,F\quad $\emptyset \subseteq \{1,2,3\}$
\wl

T\,/\,F\quad $\{1,2,3\} \subseteq \{1,2,3\}$
\wl

T\,/\,F\quad $\{1,2\} \subseteq \{1,\{2\}\}$ \hfill \textit{(Look carefully at the second set.)}
\wl

\textbf{(c)}\quad Fill in the blanks, ordering $\mathbb{N}$, $\mathbb{Z}$, $\mathbb{Q}$, $\mathbb{R}$ from ``smallest'' to ``largest'' using $\subseteq$:
\[
  \underline{\hspace{1.8cm}}
  \;\subseteq\;
  \underline{\hspace{1.8cm}}
  \;\subseteq\;
  \underline{\hspace{1.8cm}}
  \;\subseteq\;
  \underline{\hspace{1.8cm}}
\]

Circle all that apply for each number:

\medskip
\renewcommand{\arraystretch}{2.2}
\begin{tabular}{|C{2.5cm}|C{3.5cm}|C{3.5cm}|C{3.5cm}|}
\hline
Number & $\in \mathbb{N}$? & $\in \mathbb{Z}$? & $\in \mathbb{Q}$? \\
\hline
$-5$          & Yes \quad No & Yes \quad No & Yes \quad No \\
\hline
$\tfrac{1}{3}$ & Yes \quad No & Yes \quad No & Yes \quad No \\
\hline
$\sqrt{2}$     & Yes \quad No & Yes \quad No & Yes \quad No \\
\hline
$0$            & Yes \quad No & Yes \quad No & Yes \quad No \\
\hline
\end{tabular}
\renewcommand{\arraystretch}{1}

\newpage
\textbf{3.}\quad Cardinality.

\textbf{(a)}\quad Compute the cardinality of each set.

\medskip
\renewcommand{\arraystretch}{2.8}
\begin{tabular}{|C{4.5cm}|C{4.5cm}|C{4.5cm}|}
\hline
$\bigl|\{a,\, e,\, i,\, o,\, u\}\bigr|$ &
$|\emptyset|$ &
$|\{\emptyset\}|$ \\
\hline
 & & \\
\hline
\end{tabular}
\renewcommand{\arraystretch}{1}

\medskip
\textit{Note:} $\emptyset$ is a set with zero elements. $\{\emptyset\}$ is a set that \textbf{contains} the empty set --- it has exactly one element. Are they the same set? Explain.
\wls{2}

\textbf{(b)}\quad Suppose $|A| = 4$ and $|B| = 6$.
Use $|A \cup B| = |A| + |B| - |A \cap B|$.

\medskip
\hspace{1.5em}(i)\quad What is the \textbf{maximum} possible $|A \cup B|$?
When does this happen?
\wl

\hspace{1.5em}(ii)\quad What is the \textbf{minimum} possible $|A \cup B|$?
When does this happen?
\wl

\textbf{(c)}\quad $\bigstar$\;\textit{Preview.}\quad
If $|A| = m$ and $|B| = n$, give a formula for $|A \times B|$.
Explain why the formula makes sense --- don't just guess.
\wls{2}

\newpage
% =============================================================================
\section*{Part 2 --- Application}

\textbf{4.}\quad Set Operations.

Let $U = \{1,2,3,4,5,6,7,8,9,10\}$,\quad
$A = \{1,2,3,4,5,6\}$,\quad
$B = \{2,4,6,8\}$.

\textbf{(a)}\quad Complete the table.

\medskip
\renewcommand{\arraystretch}{2.4}
\begin{tabular}{|C{4cm}|C{9.5cm}|}
\hline
Expression & Result (roster notation) \\
\hline
$A \cup B$ & \\
\hline
$A \cap B$ & \\
\hline
$A \setminus B$ & \\
\hline
$B \setminus A$ & \\
\hline
$A^c$ (in $U$) & \\
\hline
$(A \cup B)^c$ (in $U$) & \\
\hline
\end{tabular}
\renewcommand{\arraystretch}{1}

\medskip
\textbf{(b)}\quad Is $A \setminus B$ the same as $B \setminus A$? Answer and explain.
\wl

\textbf{(c)}\quad Verify the formula $|A \cup B| = |A| + |B| - |A \cap B|$ using your table.
Show the arithmetic.
\wl

\textbf{(d)}\quad The reading defined complement as
``everything NOT in $A$.''
But question 4(a) says ``$A^c$ in $U$'' --- why did we need to specify $U$?
What would $A^c$ mean if we used $U' = \mathbb{N}$ as the universal set instead?
\wls{2}

\bigskip
\textbf{5.}\quad Real Life Application --- Minecraft Achievement Board.

You manage the achievement board for a Minecraft server with 30 active players.
After a big weekend event, you collect the following data:
\begin{itemize}
  \item 18 players have built a house. Call this set $H$.
  \item 14 players have defeated the Ender Dragon. Call this set $D$.
  \item 6 players have done \textbf{both}.
\end{itemize}

\textbf{(a)}\quad How many players earned at least one achievement?
\wls{2}

\textbf{(b)}\quad How many players earned \textbf{neither} achievement?
\wl

\textbf{(c)}\quad Write an expression in set notation for
``players who built a house but did NOT defeat the Ender Dragon,''
then compute the count.
\[
  \underline{\hspace{4.5cm}} \;=\; \underline{\hspace{1.2cm}} \text{ players}
\]

\textbf{(e)}\quad A third achievement is added: Nether Portal (set $P$) with $|P| = 10$.
The server rules require having a house before you can build a portal, so $P \subseteq H$.

\medskip
\hspace{1.5em}(i)\quad Can you simplify $P \cup H$? What is it equal to? Explain.
\wl

\hspace{1.5em}(ii)\quad What do you know about $P \cap H$? Can you simplify it?
\wl

\newpage
\hspace{1.5em}(iii)\quad What is the minimum possible value of $|P \cap D|$?
The maximum? (You do not have information about $P$ and $D$ specifically ---
think about what's possible.)
\wls{3}

% =============================================================================
\section*{Part 3 --- Challenge}

\textbf{6.}\quad Cartesian Products.

Let $A = \{0,\, 1,\, 2\}$ and $B = \{\text{red},\, \text{blue}\}$.

\textbf{(a)}\quad List all elements of $A \times B$.
\wls{2}

\textbf{(b)}\quad List all elements of $B \times A$.
\wls{2}

\textbf{(c)}\quad Is $A \times B = B \times A$? If not, give a specific element that is in one but not the other.
\wl

\textbf{(d)}\quad The coordinate plane is $\mathbb{R} \times \mathbb{R} = \mathbb{R}^2$.
Every point $(x,y)$ on the grid is an element of $\mathbb{R}^2$.
Write the $x$-axis as a \textit{subset} of $\mathbb{R}^2$ in set-builder notation.
\[
  x\text{-axis} = \bigl\{\;\underline{\hspace{8cm}}\;\bigr\}
\]

\textbf{(e)}\quad \textit{Preview.}\quad
Consider the function $f\colon \{1,2,3\} \to \{x,y,z\}$ defined by
$f(1) = x$, $f(2) = x$, $f(3) = z$.
Write $f$ as a subset of $\{1,2,3\} \times \{x,y,z\}$ --- just list the ordered pairs.
\[
  f = \bigl\{\;\underline{\hspace{9cm}}\;\bigr\}
\]

\bigskip
\textbf{7.}\quad \textbf{Your First Set Theory Proof.}

We are going to prove that for any sets $A$ and $B$,\quad $A \cap B \;\subseteq\; A$.

\textbf{(a)}\quad Fill in the blanks to complete the proof.

\noindent\textbf{Proof.}

\medskip
\textit{Let} $x \in A \cap B$.
\hfill \textit{(We introduce an arbitrary element of $A \cap B$.)}

\medskip
By the definition of intersection,\;
$x \in \underline{\hspace{2cm}}$\; and\; $x \in \underline{\hspace{2cm}}$.

\medskip
Therefore,\; $x \in \underline{\hspace{2cm}}$.
\hfill \textit{(This is what we needed.)}

\medskip
Since $x$ was arbitrary,\; $A \cap B \subseteq A$. Q.E.D.

\smallskip
\noindent\rule{\linewidth}{0.4pt}

\textbf{(b)}\quad Write the analogous proof that $A \cap B \subseteq B$.
(Same template --- do it yourself from scratch.)
\wls{4}

\newpage
\textbf{(c)}\quad Is $A \subseteq A \cap B$ always true?
Give a \textbf{proof} if yes, or a \textbf{counterexample} if no.
A counterexample means you find specific sets $A$ and $B$ that break it.
\wls{2}

\bigskip
\textbf{8.}\quad What does it mean for something to be \textbf{countable}? Provide a formal mathematical \textbf{definition} on what it means to be countable?
\wls{3}

% =============================================================================
\section*{Part 4 --- Notebook Artifact}

\textbf{N1.}\quad For the notes you took while doing the readings/homework (I hope you did), write at least 5 questions about any topic
today in a different color pen in your notebook.

\bigskip

\textbf{N2.}\quad If you used the internet or AI for help, describe how you used it:

\wls{3}

% =============================================================================
% =============================================================================
% =============================================================================
\end{document}
